\setcounter{section}{0}

\Needspace{5\baselineskip}
\hypertarget{ux57faux4e8eux63d0ux793aux7684ux65b9ux6cd5}{%
\section{基于提示的方法}\label{ux57faux4e8eux63d0ux793aux7684ux65b9ux6cd5}}

当下，大模型的推理能力正变得越来越重要。目前模型的推理体现为不断生成token（模型显示正在思考时也在生成token，只不过它们被包裹在特殊的标记中，因此API调用时也会计费），即在``自言自语''中进行显式的逻辑推进。以下是一些基于提示增强模型推理能力的方法。

\Needspace{5\baselineskip}
\hypertarget{ux4e00ux601dux7ef4ux94fechain-of-thought}{%
\subsection{一、思维链（Chain of Thought）}\label{ux4e00ux601dux7ef4ux94fechain-of-thought}}

要求模型逐步思考，将``大问题''拆解为``小步骤''，缩短每一步需要在嵌入空间中``跨越''的距离，进而可以通过每次``预测下一步''实现完整推理。

思维链是大模型推理技术的基石。但其局限性在于，每一步生成只靠先前已生成的内容``从前向后''推断，缺乏从``后续目标''对``先前需要的步骤''的倒推；缺乏自我纠错和反思机制，容易导致错误积累；路径单一，缺乏探索，难以覆盖复杂多样化思考路径。这三种局限性分别对应后面提到的三类方法。

\Needspace{5\baselineskip}
\hypertarget{ux4e8cux4efbux52a1ux89c4ux5212}{%
\subsection{二、任务规划}\label{ux4e8cux4efbux52a1ux89c4ux5212}}

\Needspace{4\baselineskip}
\begin{enumerate}
\def\labelenumi{\arabic{enumi}.}
\tightlist
\item
  规划-执行范式
\end{enumerate}

将推理过程显式地拆分为``规划''和``执行''两个阶段。

规划阶段：模型先不急于计算，而是根据原问题生成一份结构化的计划，列出完成任务所需的关键步骤。这相当于构建了一个全局的指导框架。

执行阶段：模型将计划中的每一步作为单独的输入，逐步生成思维链并得出结果。

优势：通过``先概览、后分步''，减少了中间环节的遗漏和计算错误，特别是在零样本（Zero-shot）场景下，能显著提升模型对未见任务的适应。

\Needspace{4\baselineskip}
\begin{enumerate}
\def\labelenumi{\arabic{enumi}.}
\setcounter{enumi}{1}
\tightlist
\item
  子任务分解
\end{enumerate}

针对难度超过示例范围的复杂问题（如多步数学推理），可以运用基于``由易到难''的递进思想。将一个复杂问题拆解为一系列子问题，并按难度递增排列，然后依次解决子问题。在解决第k个问题时，会将前k-1个问题的答案作为上下文输入。

这建立了跨步骤的显性依赖关系，利用前序结果辅助当前推理，避免逻辑断裂。

\Needspace{5\baselineskip}
\hypertarget{ux4e09ux52a8ux6001ux8c03ux6574ux4e0eux53cdux601d}{%
\subsection{三、动态调整与反思}\label{ux4e09ux52a8ux6001ux8c03ux6574ux4e0eux53cdux601d}}

如果说前面的任务规划让模型能``展望未来''，那动态调整和反思就是让模型``回看过去''，反思自己过去推理轨迹的错误并修正。

早期这通过提示词实现，后期模型通过强化学习直接学会了在推理轨迹中融入动态调整与反思过程。

\FloatBarrier
\Needspace{5\baselineskip}
\hypertarget{ux53c2ux8003ux6587ux732e-119}{%
\subsection{参考文献}\label{ux53c2ux8003ux6587ux732e-119}}

\vspace{0.25em}
\begin{itemize}
\tightlist
\item
  Wei, J., Wang, X., Schuurmans, D., et al.~(2022). \href{https://arxiv.org/abs/2201.11903}{Chain-of-Thought Prompting Elicits Reasoning in Large Language Models}. NeurIPS 2022.
\item
  Wang, L., Xu, W., Lan, Y., et al.~(2023). \href{https://aclanthology.org/2023.acl-long.147/}{Plan-and-Solve Prompting: Improving Zero-Shot Chain-of-Thought Reasoning by Large Language Models}. ACL 2023. DOI: 10.18653/v1/2023.acl-long.147.
\item
  Zhou, D., Schärli, N., Hou, L., et al.~(2023). \href{https://openreview.net/forum?id=WZH7099tgfM}{Least-to-Most Prompting Enables Complex Reasoning in Large Language Models}. ICLR 2023.
\item
  Shinn, N., Cassano, F., Gopinath, A., Narasimhan, K., \& Yao, S. (2023). \href{https://arxiv.org/abs/2303.11366}{Reflexion: Language Agents with Verbal Reinforcement Learning}. NeurIPS 2023.
\end{itemize}

\clearpage
